\documentclass[12pt]{article}
\usepackage[utf8]{inputenc}
\usepackage[T1]{fontenc}
\usepackage[margin=1in]{geometry}
\usepackage{setspace}
\usepackage{xcolor}
\usepackage[colorlinks=true, urlcolor=blue!60!black, linkcolor=blue!60!black]{hyperref}
\setlength{\parindent}{0em}
\setlength{\parskip}{0.6em}
\renewcommand{\baselinestretch}{1.0}
\begin{document}
\begin{center}
{\LARGE \textit{Dolan v. City of Tigard}, 512 U.S. 374 (1994)\par}
\vspace{0.5em}
{\large Case Brief --- Two-Page Edition\par}
\end{center}

\section{Memory Jogger}
A city that conditions a building permit on a land dedication must show rough proportionality between the dedication and the project's impact; Tigard had not (\href{https://supreme.justia.com/cases/federal/us/512/374/}{Justia}).

\section{Detailed Case Facts}
Florence Dolan sought to nearly double her plumbing and electric supply store beside Fanno Creek in Tigard, Oregon. The city required her to dedicate the parcel's floodplain as a public greenway and a 15-foot strip for a public bike path---about 10\% of her land---finding that the path ``could offset some of the traffic demand.''

\section{Procedural History}
The Oregon land-use board and courts upheld the conditions as reasonably related to the expansion's impacts.

\section{Judicial Votes}
Reversed, 5--4. Rehnquist wrote for O'Connor, Scalia, Kennedy, and Thomas; Stevens (with Blackmun and Ginsburg) and Souter dissented.

\section{Holding}
Beyond \textit{Nollan}'s essential nexus, the government must make an individualized determination that a required dedication is related in nature and extent to the development's impact. Neither condition met that standard.

\section{Analysis of Opinions}
\textbf{Majority:} Taking the strip outright would have been a per se taking, so the unconstitutional conditions doctrine applied to this adjudicative permit decision. The conditions had a nexus to flood control and traffic, but the city bore the burden of rough proportionality. It never explained why flood control required public access rather than private open space, and a path that ``could'' offset traffic was not shown to do so. \textbf{Stevens:} the Court shifted the burden to cities, ignored the benefits Dolan received, and revived \textit{Lochner}. \textbf{Souter:} the defects were really nexus failures. \emph{Later:} \textit{Koontz} (2013) and \textit{Sheetz} (2024) extended the test to money demands and legislative fees.

\section{Examples: Future Applications}
\textbf{Same-Side:} a \$60,000 traffic fee on a backyard cottage that adds about \$3,000 of traffic burden is invalid. \textbf{Opposite-Side:} a turn-lane dedication supported by a traffic study is valid. \textbf{Fence-Sitter:} a citywide school-impact fee charged to a 55-and-older community.

\section{Critique}
The burden shift may push cities to deny permits outright or to rely on broad taxes, and \textit{Sheetz} leaves unclear how much aggregation a legislative fee formula may use.

\section{Key Quotations}
\begin{quote}
``We think a term such as `rough proportionality' best encapsulates what we hold to be the requirement of the Fifth Amendment. No precise mathematical calculation is required, but the city must make some sort of individualized determination that the required dedication is related both in nature and extent to the impact of the proposed development.''
\end{quote}
\noindent --- Chief Justice Rehnquist, for the Court, 512 U.S. at 391.

\end{document}
